\thispagestyle{plain}
\begin{fullwidth}
\begin{center}
    \LARGE
    \textcolor{itesodblue}{\textbf{\mytitle}}\\
    \Large
    \vspace{0.4cm}
    \textcolor{itesomblue}{\textbf{\myauthor}}\\
    \vspace{0.9cm}
    \textcolor{itesolblue}{\textbf{Abstract}}
\end{center}
Graph-structured data describes a domain as a set of entities together with the
relations among them, a form that tabular, grid, and sequential representations
do not capture: graphs vary in size, admit no canonical ordering of their
elements, and connect each entity to a different number of others. A standard
neural network, whose input dimension and coordinate ordering are fixed, cannot
operate on such data without discarding the relational structure that carries
much of the signal. Graph Neural Networks (GNNs) are neural network
architectures designed to operate directly on graphs, respecting these
properties instead of discarding them.

This work provides a self-contained reference on GNNs for homogeneous graphs
under the message-passing framework. The foundations are established first: how
graphs represent relational data, why the multilayer perceptron (MLP) fails on
it, and the permutation invariance and equivariance that any layer defined on a
graph must satisfy. The message-passing framework is then developed as the
unifying formulation, decomposed into message, aggregation, and update
functions, together with the receptive field that successive layers induce, the
distinction between the transductive and inductive settings, and the ceiling on
expressive power.

Three representative architectures are derived as instances of the framework,
one for each standard level of graph learning: the Graph Convolutional Network (GCN) for
node classification, the Graph Auto-Encoder (GAE) for link prediction, and the
Graph Isomorphism Network (GIN) for graph classification. Each is implemented in
PyTorch Geometric on a standard benchmark --- Cora, Amazon Computers, and NCI1,
respectively --- and evaluated against an MLP baseline that predicts from node
features alone; in every task the architecture that uses the graph improves on
the baseline that does not. The limitations of the framework are examined last:
oversmoothing, oversquashing, heterophily, and the bound on expressive power.
\end{fullwidth}
